\documentclass[12pt,a4paper]{report}

\usepackage[margin=1in]{geometry}
\usepackage{xcolor}
\usepackage{sectsty}
\usepackage{fancyhdr}

% Chapter formatting
\chapterfont{\color{red}\bfseries\Huge}

% Section formatting
\sectionfont{\color{magenta}\itshape\Large}

% Subsection formatting
\subsectionfont{\color{orange}}

% Header and footer
\pagestyle{fancy}
\fancyhf{}
\fancyhead[L]{Ketan Sharma}
\fancyhead[C]{Introduction to LaTeX}
\fancyhead[R]{\today}
\fancyfoot[C]{\thepage}

\begin{document}
	
	% Title page
	\begin{titlepage}
		\centering
		
		\vspace*{2cm}
		
		{\Huge\bfseries Introduction to LaTeX\par}
		
		\vspace{1.5cm}
		
		{\Large Ketan Sharma\par}
		
		\vspace{0.5cm}
		
		$ID: YOUR_ID$
		
		\vspace{1cm}
		
		Course: LaTeX
		
		\vspace{0.5cm}
		
		Course Instructor: Instructor Name
		
		\vfill
		
		{\large \today\par}
		
	\end{titlepage}
	
	% No header/footer on first page after title page
	\thispagestyle{empty}
	
	\begin{abstract}
		This document provides a brief introduction to LaTeX and demonstrates
		the use of different document formatting features. It covers title pages,
		abstracts, chapters, sections, subsections, colors, headers, footers,
		and page formatting.
	\end{abstract}
	
	\chapter{Introduction}
	
	LaTeX is a document preparation system commonly used for writing
	technical and scientific documents. It provides powerful tools for
	formatting text, mathematical equations, tables, figures, and references.
	
	\section{What is LaTeX}
	
	LaTeX is based on the TeX typesetting system. Instead of manually
	formatting every part of a document, the user writes commands that
	describe how the document should be structured.
	
	\subsection{Advantages of LaTeX}
	
	Some advantages of LaTeX are its professional typesetting, excellent
	mathematical support, automatic numbering, and easy document
	organization.
	
	\chapter{Document Structure}
	
	A LaTeX document is divided into different structural components.
	The report class supports chapters, sections, and subsections.
	
	\section{Chapters}
	
	Chapters are the major divisions of a report document. They are created
	using the \texttt{\textbackslash chapter} command.
	
	\subsection{Sections and Subsections}
	
	Sections divide chapters into smaller topics, while subsections provide
	further organization within a section.
	
	\chapter{Formatting in LaTeX}
	
	LaTeX allows users to customize the appearance of different parts of
	a document.
	
	\section{Colors and Fonts}
	
	Packages such as \texttt{xcolor} and \texttt{sectsty} can be used to
	change the colors, sizes, and styles of headings.
	
	\subsection{Headers and Footers}
	
	The \texttt{fancyhdr} package can be used to create customized headers
	and footers. In this document, the header contains the author's name,
	course title, and current date.
	
	\chapter{Conclusion}
	
	LaTeX is a useful tool for creating well-structured and professionally
	formatted documents. Its support for chapters, sections, mathematical
	expressions, and customizable formatting makes it suitable for academic
	and technical work.
	
\end{document}